% Options for packages loaded elsewhere
\PassOptionsToPackage{unicode}{hyperref}
\PassOptionsToPackage{hyphens}{url}
\PassOptionsToPackage{dvipsnames,svgnames,x11names}{xcolor}
%
\documentclass[
  letterpaper,
  DIV=11,
  numbers=noendperiod]{scrartcl}

\usepackage{amsmath,amssymb}
\usepackage{iftex}
\ifPDFTeX
  \usepackage[T1]{fontenc}
  \usepackage[utf8]{inputenc}
  \usepackage{textcomp} % provide euro and other symbols
\else % if luatex or xetex
  \usepackage{unicode-math}
  \defaultfontfeatures{Scale=MatchLowercase}
  \defaultfontfeatures[\rmfamily]{Ligatures=TeX,Scale=1}
\fi
\usepackage{lmodern}
\ifPDFTeX\else  
    % xetex/luatex font selection
\fi
% Use upquote if available, for straight quotes in verbatim environments
\IfFileExists{upquote.sty}{\usepackage{upquote}}{}
\IfFileExists{microtype.sty}{% use microtype if available
  \usepackage[]{microtype}
  \UseMicrotypeSet[protrusion]{basicmath} % disable protrusion for tt fonts
}{}
\makeatletter
\@ifundefined{KOMAClassName}{% if non-KOMA class
  \IfFileExists{parskip.sty}{%
    \usepackage{parskip}
  }{% else
    \setlength{\parindent}{0pt}
    \setlength{\parskip}{6pt plus 2pt minus 1pt}}
}{% if KOMA class
  \KOMAoptions{parskip=half}}
\makeatother
\usepackage{xcolor}
\setlength{\emergencystretch}{3em} % prevent overfull lines
\setcounter{secnumdepth}{-\maxdimen} % remove section numbering
% Make \paragraph and \subparagraph free-standing
\makeatletter
\ifx\paragraph\undefined\else
  \let\oldparagraph\paragraph
  \renewcommand{\paragraph}{
    \@ifstar
      \xxxParagraphStar
      \xxxParagraphNoStar
  }
  \newcommand{\xxxParagraphStar}[1]{\oldparagraph*{#1}\mbox{}}
  \newcommand{\xxxParagraphNoStar}[1]{\oldparagraph{#1}\mbox{}}
\fi
\ifx\subparagraph\undefined\else
  \let\oldsubparagraph\subparagraph
  \renewcommand{\subparagraph}{
    \@ifstar
      \xxxSubParagraphStar
      \xxxSubParagraphNoStar
  }
  \newcommand{\xxxSubParagraphStar}[1]{\oldsubparagraph*{#1}\mbox{}}
  \newcommand{\xxxSubParagraphNoStar}[1]{\oldsubparagraph{#1}\mbox{}}
\fi
\makeatother

\usepackage{color}
\usepackage{fancyvrb}
\newcommand{\VerbBar}{|}
\newcommand{\VERB}{\Verb[commandchars=\\\{\}]}
\DefineVerbatimEnvironment{Highlighting}{Verbatim}{commandchars=\\\{\}}
% Add ',fontsize=\small' for more characters per line
\usepackage{framed}
\definecolor{shadecolor}{RGB}{241,243,245}
\newenvironment{Shaded}{\begin{snugshade}}{\end{snugshade}}
\newcommand{\AlertTok}[1]{\textcolor[rgb]{0.68,0.00,0.00}{#1}}
\newcommand{\AnnotationTok}[1]{\textcolor[rgb]{0.37,0.37,0.37}{#1}}
\newcommand{\AttributeTok}[1]{\textcolor[rgb]{0.40,0.45,0.13}{#1}}
\newcommand{\BaseNTok}[1]{\textcolor[rgb]{0.68,0.00,0.00}{#1}}
\newcommand{\BuiltInTok}[1]{\textcolor[rgb]{0.00,0.23,0.31}{#1}}
\newcommand{\CharTok}[1]{\textcolor[rgb]{0.13,0.47,0.30}{#1}}
\newcommand{\CommentTok}[1]{\textcolor[rgb]{0.37,0.37,0.37}{#1}}
\newcommand{\CommentVarTok}[1]{\textcolor[rgb]{0.37,0.37,0.37}{\textit{#1}}}
\newcommand{\ConstantTok}[1]{\textcolor[rgb]{0.56,0.35,0.01}{#1}}
\newcommand{\ControlFlowTok}[1]{\textcolor[rgb]{0.00,0.23,0.31}{\textbf{#1}}}
\newcommand{\DataTypeTok}[1]{\textcolor[rgb]{0.68,0.00,0.00}{#1}}
\newcommand{\DecValTok}[1]{\textcolor[rgb]{0.68,0.00,0.00}{#1}}
\newcommand{\DocumentationTok}[1]{\textcolor[rgb]{0.37,0.37,0.37}{\textit{#1}}}
\newcommand{\ErrorTok}[1]{\textcolor[rgb]{0.68,0.00,0.00}{#1}}
\newcommand{\ExtensionTok}[1]{\textcolor[rgb]{0.00,0.23,0.31}{#1}}
\newcommand{\FloatTok}[1]{\textcolor[rgb]{0.68,0.00,0.00}{#1}}
\newcommand{\FunctionTok}[1]{\textcolor[rgb]{0.28,0.35,0.67}{#1}}
\newcommand{\ImportTok}[1]{\textcolor[rgb]{0.00,0.46,0.62}{#1}}
\newcommand{\InformationTok}[1]{\textcolor[rgb]{0.37,0.37,0.37}{#1}}
\newcommand{\KeywordTok}[1]{\textcolor[rgb]{0.00,0.23,0.31}{\textbf{#1}}}
\newcommand{\NormalTok}[1]{\textcolor[rgb]{0.00,0.23,0.31}{#1}}
\newcommand{\OperatorTok}[1]{\textcolor[rgb]{0.37,0.37,0.37}{#1}}
\newcommand{\OtherTok}[1]{\textcolor[rgb]{0.00,0.23,0.31}{#1}}
\newcommand{\PreprocessorTok}[1]{\textcolor[rgb]{0.68,0.00,0.00}{#1}}
\newcommand{\RegionMarkerTok}[1]{\textcolor[rgb]{0.00,0.23,0.31}{#1}}
\newcommand{\SpecialCharTok}[1]{\textcolor[rgb]{0.37,0.37,0.37}{#1}}
\newcommand{\SpecialStringTok}[1]{\textcolor[rgb]{0.13,0.47,0.30}{#1}}
\newcommand{\StringTok}[1]{\textcolor[rgb]{0.13,0.47,0.30}{#1}}
\newcommand{\VariableTok}[1]{\textcolor[rgb]{0.07,0.07,0.07}{#1}}
\newcommand{\VerbatimStringTok}[1]{\textcolor[rgb]{0.13,0.47,0.30}{#1}}
\newcommand{\WarningTok}[1]{\textcolor[rgb]{0.37,0.37,0.37}{\textit{#1}}}

\providecommand{\tightlist}{%
  \setlength{\itemsep}{0pt}\setlength{\parskip}{0pt}}\usepackage{longtable,booktabs,array}
\usepackage{calc} % for calculating minipage widths
% Correct order of tables after \paragraph or \subparagraph
\usepackage{etoolbox}
\makeatletter
\patchcmd\longtable{\par}{\if@noskipsec\mbox{}\fi\par}{}{}
\makeatother
% Allow footnotes in longtable head/foot
\IfFileExists{footnotehyper.sty}{\usepackage{footnotehyper}}{\usepackage{footnote}}
\makesavenoteenv{longtable}
\usepackage{graphicx}
\makeatletter
\newsavebox\pandoc@box
\newcommand*\pandocbounded[1]{% scales image to fit in text height/width
  \sbox\pandoc@box{#1}%
  \Gscale@div\@tempa{\textheight}{\dimexpr\ht\pandoc@box+\dp\pandoc@box\relax}%
  \Gscale@div\@tempb{\linewidth}{\wd\pandoc@box}%
  \ifdim\@tempb\p@<\@tempa\p@\let\@tempa\@tempb\fi% select the smaller of both
  \ifdim\@tempa\p@<\p@\scalebox{\@tempa}{\usebox\pandoc@box}%
  \else\usebox{\pandoc@box}%
  \fi%
}
% Set default figure placement to htbp
\def\fps@figure{htbp}
\makeatother

\KOMAoption{captions}{tableheading}
\makeatletter
\@ifpackageloaded{caption}{}{\usepackage{caption}}
\AtBeginDocument{%
\ifdefined\contentsname
  \renewcommand*\contentsname{Table of contents}
\else
  \newcommand\contentsname{Table of contents}
\fi
\ifdefined\listfigurename
  \renewcommand*\listfigurename{List of Figures}
\else
  \newcommand\listfigurename{List of Figures}
\fi
\ifdefined\listtablename
  \renewcommand*\listtablename{List of Tables}
\else
  \newcommand\listtablename{List of Tables}
\fi
\ifdefined\figurename
  \renewcommand*\figurename{Figure}
\else
  \newcommand\figurename{Figure}
\fi
\ifdefined\tablename
  \renewcommand*\tablename{Table}
\else
  \newcommand\tablename{Table}
\fi
}
\@ifpackageloaded{float}{}{\usepackage{float}}
\floatstyle{ruled}
\@ifundefined{c@chapter}{\newfloat{codelisting}{h}{lop}}{\newfloat{codelisting}{h}{lop}[chapter]}
\floatname{codelisting}{Listing}
\newcommand*\listoflistings{\listof{codelisting}{List of Listings}}
\makeatother
\makeatletter
\makeatother
\makeatletter
\@ifpackageloaded{caption}{}{\usepackage{caption}}
\@ifpackageloaded{subcaption}{}{\usepackage{subcaption}}
\makeatother

\usepackage{bookmark}

\IfFileExists{xurl.sty}{\usepackage{xurl}}{} % add URL line breaks if available
\urlstyle{same} % disable monospaced font for URLs
\hypersetup{
  pdftitle={Vendor \& Sales Analytics},
  pdfauthor={Christian Lira Gonzalez},
  colorlinks=true,
  linkcolor={blue},
  filecolor={Maroon},
  citecolor={Blue},
  urlcolor={Blue},
  pdfcreator={LaTeX via pandoc}}


\title{Vendor \& Sales Analytics}
\author{Christian Lira Gonzalez}
\date{}

\begin{document}
\maketitle


\begin{Shaded}
\begin{Highlighting}[]
\ImportTok{import}\NormalTok{ seaborn }\ImportTok{as}\NormalTok{ sns}
\ImportTok{import}\NormalTok{ matplotlib.pyplot }\ImportTok{as}\NormalTok{ plt}
\ImportTok{import}\NormalTok{ pandas }\ImportTok{as}\NormalTok{ pd}
\ImportTok{from}\NormalTok{ pathlib }\ImportTok{import}\NormalTok{ Path}
\ImportTok{import}\NormalTok{ pyodbc}
\end{Highlighting}
\end{Shaded}

\begin{Shaded}
\begin{Highlighting}[]
\NormalTok{conn }\OperatorTok{=}\NormalTok{ pyodbc.}\ExtensionTok{connect}\NormalTok{(}
    \StringTok{"DRIVER=\{ODBC Driver 17 for SQL Server\};"}
    \StringTok{"SERVER=localhost;"}
    \StringTok{"DATABASE=Northwind;"}
    \StringTok{"Trusted\_Connection=yes;"}
\NormalTok{)}

\CommentTok{\# Import supplier\_ranking query from SQL files in dir}
\NormalTok{cohort\_sql }\OperatorTok{=}\NormalTok{ Path(}\StringTok{"cohort\_analysis.sql"}\NormalTok{).read\_text()}
\NormalTok{cohort\_sql }\OperatorTok{=}\NormalTok{ pd.read\_sql(cohort\_sql, conn)}

\NormalTok{conn.close()}
\end{Highlighting}
\end{Shaded}

\subsection{Cohort Analysis SQL Logic}\label{cohort-analysis-sql-logic}

The SQL query above performs a customer cohort retention analysis with
the following steps:

\begin{enumerate}
\def\labelenumi{\arabic{enumi}.}
\item
  \textbf{FirstPurchase CTE} -- Identifies each customer's first
  purchase date and truncates it to the first day of that month (the
  cohort month).
\item
  \textbf{CustomerActivity CTE} -- Joins the original orders to each
  customer's cohort, calculating the number of months elapsed between
  the cohort month and each subsequent order month.
\item
  \textbf{CohortSize CTE} -- Counts the total number of unique customers
  in each cohort month.
\item
  \textbf{Final SELECT} -- For each cohort month and each month elapsed
  (0 to 6 months), it calculates:

  \begin{itemize}
  \tightlist
  \item
    \texttt{CohortSize} -- Total customers in that cohort
  \item
    \texttt{ActiveCustomers} -- Customers who made a purchase in that
    period
  \item
    \texttt{RetentionRate} -- Percentage of customers who remained
    active (ActiveCustomers / CohortSize × 100)
  \end{itemize}
\end{enumerate}

The \texttt{WHERE\ ca.MonthsElapsed\ BETWEEN\ 0\ AND\ 6} limits the
analysis to the first 6 months after each customer's first purchase,
which is typically the most telling period for retention patterns.

\begin{Shaded}
\begin{Highlighting}[]
\OperatorTok{{-}{-}}\NormalTok{ USE Northwind}\OperatorTok{;}
\OperatorTok{{-}{-}}\NormalTok{ GO}

\NormalTok{WITH FirstPurchase AS (}
\NormalTok{    SELECT}
\NormalTok{        CustomerID,}
\NormalTok{        DATEFROMPARTS(YEAR(MIN(OrderDate)), MONTH(MIN(OrderDate)), }\DecValTok{1}\NormalTok{) AS CohortMonth}
\NormalTok{    FROM Orders}
\NormalTok{    WHERE OrderDate IS NOT NULL}
\NormalTok{    GROUP BY CustomerID}
\NormalTok{),}
\NormalTok{CustomerActivity AS (}
\NormalTok{    SELECT}
\NormalTok{        fp.CustomerID,}
\NormalTok{        fp.CohortMonth,}
\NormalTok{        DATEFROMPARTS(YEAR(o.OrderDate), MONTH(o.OrderDate), }\DecValTok{1}\NormalTok{) AS OrderMonth,}
\NormalTok{        DATEDIFF(MONTH, fp.CohortMonth,}
\NormalTok{            DATEFROMPARTS(YEAR(o.OrderDate), MONTH(o.OrderDate), }\DecValTok{1}\NormalTok{)) AS MonthsElapsed}
\NormalTok{    FROM FirstPurchase fp}
\NormalTok{    JOIN Orders o ON fp.CustomerID }\OperatorTok{=}\NormalTok{ o.CustomerID}
\NormalTok{    WHERE o.OrderDate IS NOT NULL}
\NormalTok{),}
\NormalTok{CohortSize AS (}
\NormalTok{    SELECT CohortMonth, COUNT(DISTINCT CustomerID) AS TotalCustomers}
\NormalTok{    FROM FirstPurchase}
\NormalTok{    GROUP BY CohortMonth}
\NormalTok{)}
\NormalTok{SELECT}
\NormalTok{    ca.CohortMonth,}
\NormalTok{    cs.TotalCustomers                                           AS CohortSize,}
\NormalTok{    ca.MonthsElapsed,}
\NormalTok{    COUNT(DISTINCT ca.CustomerID)                              AS ActiveCustomers,}
\NormalTok{    ROUND(}\FloatTok{100.0} \OperatorTok{*}\NormalTok{ COUNT(DISTINCT ca.CustomerID)}
        \OperatorTok{/}\NormalTok{ cs.TotalCustomers,}\DecValTok{2}\NormalTok{)                               AS RetentionRate}
\NormalTok{FROM CustomerActivity ca}
\NormalTok{JOIN CohortSize cs ON ca.CohortMonth }\OperatorTok{=}\NormalTok{ cs.CohortMonth}
\NormalTok{WHERE ca.MonthsElapsed BETWEEN }\DecValTok{0}\NormalTok{ AND }\DecValTok{6}
\NormalTok{GROUP BY ca.CohortMonth, cs.TotalCustomers, ca.MonthsElapsed}
\NormalTok{ORDER BY ca.CohortMonth, ca.MonthsElapsed}\OperatorTok{;}
\end{Highlighting}
\end{Shaded}

\begin{Shaded}
\begin{Highlighting}[]
\BuiltInTok{print}\NormalTok{(cohort\_sql.head(}\DecValTok{10}\NormalTok{))}
\end{Highlighting}
\end{Shaded}

\begin{verbatim}
  CohortMonth  CohortSize  MonthsElapsed  ActiveCustomers  RetentionRate
0  1996-07-01          20              0               20         100.00
1  1996-07-01          20              1                4          20.00
2  1996-07-01          20              2                4          20.00
3  1996-07-01          20              3                4          20.00
4  1996-07-01          20              4                4          20.00
5  1996-07-01          20              5                6          30.00
6  1996-07-01          20              6                7          35.00
7  1996-08-01          14              0               14         100.00
8  1996-08-01          14              1                6          42.86
9  1996-08-01          14              2                4          28.57
\end{verbatim}

\subsection{Raw Retention Data
Interpretation}\label{raw-retention-data-interpretation}

The table above shows the following structure:

\begin{itemize}
\tightlist
\item
  \textbf{CohortMonth} -- The month of the customer's first purchase
  (e.g., July 1996, August 1996)
\item
  \textbf{CohortSize} -- Number of new customers acquired that month
\item
  \textbf{MonthsElapsed} -- How many months after first purchase (0 =
  first purchase month)
\item
  \textbf{ActiveCustomers} -- Number of customers from that cohort who
  purchased in that period
\item
  \textbf{RetentionRate} -- Percentage retention (ActiveCustomers /
  CohortSize × 100)
\end{itemize}

\textbf{Key observations from the raw data:} - July 1996 cohort had 20
new customers. Retention dropped sharply from 100\% (Month 0) to 20\%
(Month 1), then gradually recovered to 35\% by Month 6. - August 1996
cohort had 14 new customers with a stronger initial retention of 42.86\%
at Month 1, though it declined to 28.57\% by Month 2.

\begin{Shaded}
\begin{Highlighting}[]
\ImportTok{import}\NormalTok{ seaborn }\ImportTok{as}\NormalTok{ sns}
\ImportTok{import}\NormalTok{ matplotlib.pyplot }\ImportTok{as}\NormalTok{ plt}
\ImportTok{import}\NormalTok{ pandas }\ImportTok{as}\NormalTok{ pd}

\NormalTok{df }\OperatorTok{=}\NormalTok{ cohort\_sql }\CommentTok{\# the T{-}SQL above}

\CommentTok{\# Pivot: rows = cohort month, columns = months elapsed, values = retention rate}
\NormalTok{cohort\_pivot }\OperatorTok{=}\NormalTok{ df.pivot\_table(}
\NormalTok{    index}\OperatorTok{=}\StringTok{\textquotesingle{}CohortMonth\textquotesingle{}}\NormalTok{,}
\NormalTok{    columns}\OperatorTok{=}\StringTok{\textquotesingle{}MonthsElapsed\textquotesingle{}}\NormalTok{,}
\NormalTok{    values}\OperatorTok{=}\StringTok{\textquotesingle{}RetentionRate\textquotesingle{}}
\NormalTok{)}

\CommentTok{\# Format cohort month labels for readability}
\NormalTok{cohort\_pivot.index }\OperatorTok{=}\NormalTok{ pd.to\_datetime(cohort\_pivot.index).strftime(}\StringTok{\textquotesingle{}\%Y{-}\%m\textquotesingle{}}\NormalTok{)}

\CommentTok{\# Plot heatmap}
\NormalTok{plt.figure(figsize}\OperatorTok{=}\NormalTok{(}\DecValTok{12}\NormalTok{, }\DecValTok{8}\NormalTok{))}
\NormalTok{sns.heatmap(}
\NormalTok{    cohort\_pivot,}
\NormalTok{    annot}\OperatorTok{=}\VariableTok{True}\NormalTok{,}
\NormalTok{    fmt}\OperatorTok{=}\StringTok{\textquotesingle{}.0f\textquotesingle{}}\NormalTok{,}
\NormalTok{    cmap}\OperatorTok{=}\StringTok{\textquotesingle{}Blues\textquotesingle{}}\NormalTok{,}
\NormalTok{    linewidths}\OperatorTok{=}\FloatTok{0.5}\NormalTok{,}
\NormalTok{    vmin}\OperatorTok{=}\DecValTok{0}\NormalTok{, vmax}\OperatorTok{=}\DecValTok{100}\NormalTok{,}
\NormalTok{    cbar\_kws}\OperatorTok{=}\NormalTok{\{}\StringTok{\textquotesingle{}label\textquotesingle{}}\NormalTok{: }\StringTok{\textquotesingle{}Retention Rate (\%)\textquotesingle{}}\NormalTok{\}}
\NormalTok{)}
\NormalTok{plt.title(}\StringTok{\textquotesingle{}Customer Cohort Retention Rate (\%)\textquotesingle{}}\NormalTok{, fontsize}\OperatorTok{=}\DecValTok{14}\NormalTok{, pad}\OperatorTok{=}\DecValTok{15}\NormalTok{)}
\NormalTok{plt.xlabel(}\StringTok{\textquotesingle{}Months Since First Purchase\textquotesingle{}}\NormalTok{)}
\NormalTok{plt.ylabel(}\StringTok{\textquotesingle{}Cohort (First Purchase Month)\textquotesingle{}}\NormalTok{)}
\NormalTok{plt.tight\_layout()}
\NormalTok{plt.savefig(}\StringTok{\textquotesingle{}cohort\_heatmap.png\textquotesingle{}}\NormalTok{, dpi}\OperatorTok{=}\DecValTok{150}\NormalTok{, bbox\_inches}\OperatorTok{=}\StringTok{\textquotesingle{}tight\textquotesingle{}}\NormalTok{)}
\NormalTok{plt.show()}
\end{Highlighting}
\end{Shaded}

\pandocbounded{\includegraphics[keepaspectratio]{CohortAnalysis_files/figure-pdf/cell-6-output-1.png}}

\subsection{Cohort Retention Heatmap
Analysis}\label{cohort-retention-heatmap-analysis}

The heatmap above visualizes customer retention rates across different
cohorts over time.

\subsubsection{How to Read This Chart:}\label{how-to-read-this-chart}

\begin{itemize}
\tightlist
\item
  \textbf{Rows} represent customer cohorts (the month they made their
  first purchase)
\item
  \textbf{Columns} represent months since the first purchase (0 to 6)
\item
  \textbf{Color intensity} (darker = higher retention percentage)
\item
  \textbf{Numbers} show the exact retention rate for each cohort at each
  interval
\end{itemize}

\subsubsection{Key Business Insights:}\label{key-business-insights}

\begin{enumerate}
\def\labelenumi{\arabic{enumi}.}
\item
  \textbf{Initial Drop-off is Severe} -- All cohorts show a dramatic
  decline from Month 0 (100\%) to Month 1 (20--43\%). This indicates
  that while initial acquisition is successful, the first post-purchase
  month is critical for customer churn.
\item
  \textbf{July 1996 Cohort Performance} -- This cohort (20 customers)
  started strong but dropped to 20\% by Month 1. However, retention
  improved to 30\% by Month 5 and 35\% by Month 6, suggesting that
  re-engagement efforts or seasonal factors may have brought customers
  back.
\item
  \textbf{August 1996 Cohort} -- Though smaller (14 customers), this
  cohort showed better initial retention (42.86\%) at Month 1,
  indicating potentially higher-quality acquisitions or a more engaging
  onboarding experience.
\item
  \textbf{Long-term Retention Trends} -- Cohorts tend to stabilize
  between 20--35\% after Month 2, suggesting that customers who remain
  beyond the second month are more likely to become loyal, repeat
  purchasers.
\end{enumerate}

\subsubsection{Recommended Actions:}\label{recommended-actions}

\begin{itemize}
\tightlist
\item
  Investigate what changed between July and August 1996 that improved
  Month 1 retention.
\item
  Implement targeted win-back campaigns for customers who haven't
  purchased since their first month.
\item
  Analyze the behavior of customers who remain active at Month 6
  vs.~those who churn early to identify key retention drivers.
\end{itemize}




\end{document}
